% !TeX program = xelatex
% Figure 5. Build with XeLaTeX. DATA ancestry and external references, not a common time axis.
\documentclass[tikz,border=6pt]{standalone}
\usepackage{lmodern}
\usepackage{fontspec,xeCJK}
\setmainfont[BoldFont=texgyreheros-bold.otf]{texgyreheros-regular.otf}
\setCJKmainfont[BoldFont=FandolHei-Bold.otf]{FandolHei-Regular.otf}
\usepackage{amsmath,amssymb}
\usetikzlibrary{arrows.meta,calc,positioning,backgrounds}
\definecolor{ink}{HTML}{263445}
\definecolor{muted}{HTML}{586779}
\definecolor{line}{HTML}{6B7C8F}
\definecolor{previous}{HTML}{EDE9FE}
\definecolor{regionA}{HTML}{DBEAFE}
\definecolor{regularB}{HTML}{FFEDD5}
\definecolor{forceB}{HTML}{FCE7F3}
\definecolor{green}{HTML}{15803D}
\definecolor{rose}{HTML}{BE185D}
\definecolor{blue}{HTML}{2563EB}
\definecolor{orange}{HTML}{C2410C}
\begin{document}
\begin{tikzpicture}[
  x=1cm,y=1cm,
  font=\fontsize{10}{13}\selectfont,text=ink,
  block/.style={draw=line,rounded corners=2pt,align=center,
    inner sep=4pt,text width=2.35cm,minimum height=1.05cm,line width=.65pt},
  ext/.style={block,fill=gray!4,minimum height=.95cm},
  kept/.style={draw=green,line width=1.65pt},
  chain/.style={-{Stealth[length=2.1mm]},line width=.9pt,draw=line},
  ref/.style={-{Stealth[length=1.8mm]},densely dashed,line width=.65pt,draw=muted},
  forced/.style={-{Stealth[length=2.6mm]},line width=1.5pt,draw=rose},
  tag/.style={fill=white,inner sep=2pt,align=center,
    font=\fontsize{9.5}{12}\selectfont},
  rowlabel/.style={anchor=east,align=right,text=muted,font=\bfseries\fontsize{10}{13}\selectfont},
  card/.style={draw=line!65,rounded corners=3pt,align=left,
    inner sep=7pt,text width=9.45cm,minimum height=2.2cm,
    font=\fontsize{10}{14}\selectfont}
]

% Title and color key.
\node[font=\bfseries\fontsize{14}{18}\selectfont] at (10.2,14.0)
  {两条 BFT 分支，共享截至 108 的全局前缀};
\node[anchor=west,font=\fontsize{9}{12}\selectfont,text=muted] at (.85,13.28)
  {填色区分 session；绿色粗框表示已进入 $G$。块内的 $(\beta,r)$ 是其生产 session。};
\node[draw=green,line width=1.5pt,fill=white,minimum width=.32cm,minimum height=.20cm,inner sep=0] at (17.2,13.28) {};
\node[anchor=west,font=\fontsize{9}{12}\selectfont,text=green] at (17.48,13.28) {$G$ 中的块保留};

% L1 chain.
\node[rowlabel] at (.55,11.95) {L1 链};
\node[ext] (l2500) at (2.3,11.95) {高度 2500\\$\mathcal C_7$ 已激活};
\node[ext] (l2600) at (8.7,11.95) {高度 2600\\仍对应 $\mathcal C_7$};
\node[ext] (l3000) at (15.3,11.95) {高度 3000\\$\mathcal C_8$ 激活};
\node[ext,text width=2.8cm] (l3100) at (19.0,11.95)
  {3100：公布 $U_{41}$\\新表 $S_8$；$m_8=8$};
\draw[chain] (l2500) -- (l2600);
\draw[chain] (l2600) -- (l3000);
\draw[chain] (l3000) -- (l3100);

% Shared global prefix and the old regular suffix.
\begin{scope}[on background layer]
  \fill[green!3,rounded corners=4pt] (.87,8.50) rectangle (10.15,10.75);
  \draw[draw=green!45,rounded corners=4pt] (.87,8.50) rectangle (10.15,10.75);
  \draw[draw=muted!65,densely dashed,rounded corners=4pt] (10.70,8.50) rectangle (20.65,10.75);
\end{scope}
\node[rowlabel] at (.55,9.30) {BFT 链\\共同前缀};
\node[font=\bfseries\fontsize{10}{13}\selectfont,text=green] at (5.5,10.50)
  {$G$：激活 $J_{501}$ 前已全局确认到 108};
\node[font=\bfseries\fontsize{10}{13}\selectfont,text=orange] at (15.65,10.50)
  {Regular 延续：$\beta=7$，始终使用 $S_7$};

\node[block,fill=previous,kept] (b100) at (2.3,9.30)
  {\textbf{DATA 100}\\$(7,19)$\\Regular 参考};
\node[block,fill=regionA,kept] (b101) at (5.5,9.30)
  {\textbf{101--107}\\A：$(7,20)$};
\node[block,fill=regionA,kept] (b108) at (8.7,9.30)
  {\textbf{DATA 108}\\A：$(7,20)$\\参考与起点};
\node[block,fill=regionA] (b109) at (12.1,9.30)
  {\textbf{旧 DATA 109}\\A：$(7,20)$};
\node[block,fill=regionA] (b110) at (15.3,9.30)
  {\textbf{旧 DATA 110}\\A：$(7,20)$\\正常接棒点};
\node[block,fill=regularB,text width=2.8cm] (b111) at (19.0,9.30)
  {\textbf{旧 111--115}\\B：$(7,21)$};
\draw[chain] (b100) -- (b101);
\draw[chain] (b101) -- (b108);
\draw[chain,draw=blue] (b108) -- (b109);
\draw[chain,draw=blue] (b109) -- (b110);
\draw[chain,draw=orange] (b110) -- (b111);
\node[tag,text=muted] at (15.65,8.15)
  {旧 109--115 未进入 $G$；有效 Force 后，这段后缀不再被接受};

% The real fork: the new path starts from the same b108 node.
\node[font=\bfseries\fontsize{10}{13}\selectfont,text=rose] at (15.65,7.40)
  {Force 从 108 分叉：$\beta=8$，安装并使用 $S_8$};
\node[block,fill=forceB,draw=rose] (f109) at (12.1,6.50)
  {\textbf{新 DATA 109}\\B：$(8,21)$};
\node[block,fill=forceB,draw=rose] (f110) at (15.3,6.50)
  {\textbf{新 DATA 110}\\B：$(8,21)$};
\node[block,fill=forceB,draw=rose,text width=2.8cm] (f111) at (19.0,6.50)
  {\textbf{新 DATA 111\ldots}\\B：$(8,21)$};
\coordinate (fork) at (10.30,9.30);
\draw[draw=green,line width=1.1pt] (b108.east) -- (fork);
\fill[green] (fork) circle (1.7pt);
\draw[forced] (fork) -- (10.30,6.50) -- (f109.west);
\draw[forced] (f109) -- (f110);
\draw[forced] (f110) -- (f111);

% External references from the actual selected BFT blocks.
\draw[ref] (b100.north) -- node[tag,right,pos=.75] {$\lambda(100)=2500$} (l2500.south);
\draw[ref] (b108.north) -- node[tag,right,pos=.75] {$\lambda(108)=2600$} (l2600.south);

% Espresso chain and the activation boundary.
\node[rowlabel] at (.55,3.80) {Espresso 链};
\node[ext] (j490) at (2.3,3.80) {$J_{490}$\\100 的 $\xi$ 引用};
\node[ext] (j495) at (5.5,3.80) {$J_{495}$\\108 的 $\xi$ 引用};
\node[ext] (j500) at (8.7,3.80) {$J_{500}$：时间 119\\验收后 $G$ 到 108};
\node[ext,draw=rose,fill=forceB!55,line width=1.35pt] (j501) at (12.1,3.80)
  {$J_{501}$：时间 122\\$U_{41}$ 生效\\$\beta:7\to8$，安装 $S_8$};
\node[ext] (j502) at (15.3,3.80) {$J_{502}$\\后续验收新分支};
\draw[chain] (j490) -- (j495);
\draw[chain] (j495) -- (j500);
\draw[chain] (j500) -- (j501);
\draw[chain] (j501) -- (j502);
\draw[ref] (b100.south) -- node[tag,left,pos=.67] {$\xi(100)$} (j490.north);
\draw[ref] (b108.south) -- node[tag,left,pos=.57] {$\xi(108)$} (j495.north);
\draw[ref,draw=rose,line width=.9pt] (j501.north) -- node[tag,right] {启用新分支\\与新状态表} (f109.south);

% Fixed inputs and explicit region selection, one compact card per rotation.
\node[card,fill=orange!4,draw=orange!70,anchor=north west] at (.88,2.35) {
  \textbf{Regular B：session $(7,21)$}\\
  固定：权重快照 $\mathcal C_7$；seed $=100$；状态表 $S_7$。\\
  $S_7$：A ON，B ON，C OFF（更早由 $U_{40}$ 固定）。\\
  顺序 C $\to$ B $\to$ A：跳过 C；B 在 $\mathcal C_7$ 中权重 $>0$，选 B。
};
\node[card,fill=rose!4,draw=rose!70,anchor=north west] at (11.02,2.35) {
  \textbf{Force B：session $(8,21)$}\\
  固定：权重快照 $\mathcal C_8$（$m_8=8$）；seed $=108$；状态表 $S_8$。\\
  $S_8$：A OFF，B ON，C OFF（随新 $\beta$ 一起安装）。\\
  启动 A $\to$ B $\to$ C：跳过 A；B 在 $\mathcal C_8$ 中权重 $>0$，选 B。
};
\node[anchor=north west,font=\fontsize{9}{12}\selectfont,text=muted] at (.88,-.05)
  {三条链各有自己的高度，横向位置不表示同时发生。$S$ 的下标是分支版本，$\mathcal C$ 的下标是快照编号。};
\end{tikzpicture}
\end{document}
